\scanpage{147}
De plus, le foncteur $M*M'$ est bi-additif en chaque
argument. Il y a un élément unité pour $*$ :
$K\in{\unsure{\mathcal A|A}}^{\unsure{\mathfrak S|\mathcal S}_0}$, \deleted{\illegible{1}} [où par définition $\mathfrak S_0=\{e\}$, et $u_{0,p,\ldots}=u_{p,\ldots}$]

Passons aux groupes de classes (pour \unsure{le splittage|les splittings},
pour fixer les idées), posons
\[
\deleted{\unsure{K(G)}}\quad A(G)=K(\mathcal A^G) \tag*{(22)}
\]
$A(G)$ est un \deleted{\unsure{faisceau}} \uline{anneau de Chern}
\uline{\unsure{involutif|commutatif}} \uline{augmenté} [de plus, il est de dimension
finie, et même de rang fini sur les entiers, mais
c'est fortuit], $A(G)$ est un foncteur
\deleted{\uline{covariant}} contravariant en $G$, $u:G\longrightarrow G'$ définit
\[
A(u):A(G')\longrightarrow A(G) \tag*{(\unsure{23})}
\]
% slanted
\margin{$A(u)$ est injectif si $u$ est surjectif (grâce au th. de Jordan-Hölder)}
et $A(1_G)=1_{A(G)}$, $A(uv)=A(v)A(u)$. Cela permet de
définir le produit extérieur $\alpha\in A(G)$, $\beta\in A(G')$
\[
\alpha\square\beta=\deleted{\unsure{\alpha}}\ A(p_1)(\alpha).\,A(p_2)(\beta)\in A(G\times G') \tag*{(\unsure{24})}
\]
D'ailleurs, la seule connaissance de la structure de
groupe sur les $A(G)$, et de l'opération externe
$\alpha\square\beta$, redonne les structures d'anneau \deleted{\illegible{1}} des $A(G)$.
[Mais pas leur structure de $\lambda$-anneau pour l'instant).

Il y a aussi les foncteurs \inserted{\uline{additifs}}, pour $u:G\longrightarrow G'$
\unsure{injectifs}
\[
I(u):A(G)\longrightarrow A(G') \tag*{(25)}
\]
(mais ce foncteur n'\uline{est pas multiplicatif})

On a
\[
\begin{aligned}
A(u)I(u)&=\text{multiplication par }[G':G]\\
I(u)A(u)&=\text{multiplication par }[G':G]
\end{aligned}\tag*{(26)}
\]
% archive sheet number: 146
% printed form (vertical paper watermark)
\unsure{GENTINE} \tikz[baseline=-.2ex,scale=.17]{\draw (0,0)--(1.6,0)--(1.8,.33)--(.5,.33)--(.32,.13)--(.92,.13);} TYPEWRITER \tikz[baseline=-.4ex,scale=.15]{\foreach \a in {0,45,...,315}{\draw (\a:.52) circle (.22);} \draw (0,0) circle (.21);} \unsure{ARGENT}
